\chapter{Conception du Data Warehouse et modélisation}

\section{Constellation retenue}
Le modèle physique recommandé est une constellation minimale :
\begin{figure}[H]
\centering
\resizebox{0.94\textwidth}{!}{%
\begin{tikzpicture}[node distance=1.2cm, every node/.style={font=\small}, dim/.style={draw=gold, fill=white, rounded corners, align=center, minimum width=3.2cm, minimum height=0.8cm}, fact/.style={draw=navy, fill=lightnavy, rounded corners, align=center, minimum width=4.3cm, minimum height=1.2cm}, rel/.style={-, thick, color=navy}]
\node[fact] (obs) {FactSensorObservation\\grain : observation source};
\node[fact, below=2.5cm of obs] (fail) {FactFailurePeriod\\grain : rapport externe};
\node[dim, left=2.7cm of obs] (date) {DimDate\\DateKey};
\node[dim, right=2.7cm of obs] (time) {DimTime\\TimeKey};
\draw[rel] (date) -- (obs);
\draw[rel] (time) -- (obs);
\draw[rel] (date.south) |- node[left,align=center]{StartDateKey\\actif} (fail.west);
\draw[rel,dashed] (date.south) |- node[right,align=center]{EndDateKey\\inactif} (fail.east);
\end{tikzpicture}
}%
\caption{Constellation dimensionnelle du projet}
\label{fig:dw}
\sourcefigure{élaboration propre à partir du modèle dimensionnel validé}
\end{figure}

\section{FactSensorObservation}
Cette table conserve ObservationKey, SourceRowKey, DateKey, TimeKey, ObservationTimestamp, FailurePeriodKey, ObservationCount et les quinze signaux. FailurePeriodKey est une colonne dérivée de contexte ; elle ne change pas le grain. ObservationCount vaut 1 par ligne et permet une mesure additive. Les signaux analogiques sont résumés par moyenne, minimum ou maximum selon le besoin ; les signaux 0/1 sont des états et non des pannes.

\section{FactFailurePeriod}
Cette table contient FailurePeriodKey, FailureReportCode, les bornes complètes, les clés de date, FailureType, Severity, MaintenanceNoteRaw, SourceReference, DurationSeconds et FailurePeriodCount. Elle reste à quatre lignes. La relation StartDateKey vers DimDate est active. EndDateKey joue un rôle de date secondaire et la relation est inactive dans le modèle recommandé.

\section{Dimensions conformes}
DimDate contient DateKey, CalendarDate, Year, QuarterNumber, MonthNumber, MonthName et DayOfMonth. DimTime contient 86\,400 secondes de la journée et les attributs TimeLabel, Hour, Minute, Second et TimeBucket. La valeur exacte reste dans ObservationTimestamp ; DimTime sert aux regroupements horaires.

\section{Relations et intégrité}
Les relations sont un-à-plusieurs, à sens de filtre simple, depuis les dimensions vers les faits. Aucun lien direct entre FactSensorObservation et FactFailurePeriod n'est nécessaire. L'association incidente est faite par comparaison temporelle dans Power Query ou Python. Cette stratégie évite une relation ambiguë et évite de multiplier les périodes.

\section{Choix et compromis}
Un format long aurait rendu naturelle une DimSignal, mais aurait théoriquement produit 22\,754\,220 lignes. Le format large conserve le grain source et limite la complexité pédagogique. Un fait agrégé de cinq minutes pourrait améliorer les tendances, mais il ne doit pas remplacer le fait original car la cadence comporte des trous et des transitions courtes.
